% GENERATED FILE. Edit canonical lessons, then run npm run generate:books.
% canonical-source-hash: fnv1a64:1c69fc4047f3f73a
% canonical-lessons: TA-C310-keliccittiram, TA-C310-cuvarotti, TA-C310-avanappatam, TA-C310-cittari, TA-C310-utakam

\chapter{Pictures and Posters}
\label{ch:ta-pictures-and-posters}
\hlchaptermodality{310}

\begin{chapteropening}
\textbf{By the end of this chapter:} \emph{I can talk about cartoons, posters and films about real life.}

Complete the last lesson of chapter 310: I can talk about cartoons, posters and films about real life.
\end{chapteropening}

\section[\texorpdfstring{\ta{கேலிச்சித்திரம்} (kēliccittiram)}{கேலிச்சித்திரம் (kēliccittiram)}]{\ta{கேலிச்சித்திரம்} (kēliccittiram) — a cartoon}

\label{lesson:TA-C310-keliccittiram}



\begin{quote}
Before the new one: say the Tamil for to state, then the Tamil for to hear of.
\end{quote}

\subsection*{You'll want to know: \ta{கேலிச்சித்திரம்}}
\textbf{\ta{கேலிச்சித்திரம்}} — \emph{kēliccittiram} — \textquotedblleft{}a cartoon\textquotedblright{}.

\begin{quote}
\ta{நாளிதழில்} \ta{ஒரு} \ta{வேடிக்கையான} \ta{கேலிச்சித்திரம்} \ta{வந்தது}.
\end{quote}

\begin{quote}
\emph{nāḷitaḻil oru vēṭikkaiyāṉa kēliccittiram vantatu.}
\end{quote}

\begin{quote}
A funny cartoon appeared in the newspaper.
\end{quote}

\ta{கேலி}, mockery + \ta{சித்திரம்}, picture.

\begin{cousinweb}
\textbf{\ta{கேலிச்சித்திரம்}} (\emph{kēliccittiram}) joins \ta{கேலி}, \textquotedblleft{}mockery\textquotedblright{}, usually traced to Sanskrit \emph{keli}, \textquotedblleft{}play, sport\textquotedblright{} (Monier-Williams), and \ta{சித்திரம்}, \textquotedblleft{}a picture\textquotedblright{}, from Sanskrit \emph{citra}, \textquotedblleft{}bright, variegated\textquotedblright{} (CDIAL 4803).
\end{cousinweb}

\begin{grammarlens}[title={What you've built}]
One more word for telling what happened.
\end{grammarlens}

\subsection*{Your turn}
\begin{itemize}
\raggedright
  \item \emph{Say it:} \emph{kēliccittiram}
  \item \emph{Say it:} \emph{kēliccittiram}, once more
  \item \emph{Say it:} \emph{nāḷitaḻil oru vēṭikkaiyāṉa kēliccittiram vantatu.} — A funny cartoon appeared in the newspaper.
  \item \emph{Recall:} say the Tamil for to state, then the Tamil for to hear of, then say \emph{kēliccittiram} again
\end{itemize}

\begin{tcolorbox}[breakable,colback=teal!4,colframe=teal!35!black,title={Before you move on}]
What does \ta{கேலிச்சித்திரம்} mean? (\textquotedblleft{}A cartoon\textquotedblright{}.) Say it once more.
\end{tcolorbox}
\section[\texorpdfstring{\ta{சுவரொட்டி} (cuvaroṭṭi)}{சுவரொட்டி (cuvaroṭṭi)}]{\ta{சுவரொட்டி} (cuvaroṭṭi) — a poster}

\label{lesson:TA-C310-cuvarotti}



\begin{quote}
Before the new one: say the Tamil for to hear of, then the Tamil for a cartoon.
\end{quote}

\subsection*{You'll want to know: \ta{சுவரொட்டி}}
\textbf{\ta{சுவரொட்டி}} — \emph{cuvaroṭṭi} — \textquotedblleft{}a poster\textquotedblright{}.

\begin{quote}
\ta{சுவர்களில்} \ta{புதிய} \ta{படத்தின்} \ta{சுவரொட்டிகள்} \ta{இருந்தன}.
\end{quote}

\begin{quote}
\emph{cuvarkaḷil putiya paṭattiṉ cuvaroṭṭikaḷ iruntaṉa.}
\end{quote}

\begin{quote}
There were posters for the new film on the walls.
\end{quote}

Plural: \ta{சுவரொட்டிகள்}, \emph{cuvaroṭṭikaḷ}.

\ta{சுவர்}, wall + \ta{ஒட்டி}, a thing stuck on.

\begin{cousinweb}
\textbf{\ta{சுவரொட்டி}} (\emph{cuvaroṭṭi}) is a Tamil compound built from \ta{சுவர்}, \textquotedblleft{}wall\textquotedblright{} (DEDR 1975), and \ta{ஒட்டு}, \textquotedblleft{}to stick, glue on\textquotedblright{} (DEDR 958).
\end{cousinweb}

\begin{grammarlens}[title={What you've built}]
One more word for telling what happened.
\end{grammarlens}

\subsection*{Your turn}
\begin{itemize}
\raggedright
  \item \emph{Say it:} \emph{cuvaroṭṭi}
  \item \emph{Say it:} \emph{cuvaroṭṭi}, once more
  \item \emph{Say it:} \emph{cuvarkaḷil putiya paṭattiṉ cuvaroṭṭikaḷ iruntaṉa.} — There were posters for the new film on the walls.
  \item \emph{Recall:} say the Tamil for to hear of, then the Tamil for a cartoon, then say \emph{cuvaroṭṭi} again
\end{itemize}

\begin{tcolorbox}[breakable,colback=teal!4,colframe=teal!35!black,title={Before you move on}]
What does \ta{சுவரொட்டி} mean? (\textquotedblleft{}A poster\textquotedblright{}.) Say it once more.
\end{tcolorbox}
\section[\texorpdfstring{\ta{ஆவணப்படம்} (āvaṇappaṭam)}{ஆவணப்படம் (āvaṇappaṭam)}]{\ta{ஆவணப்படம்} (āvaṇappaṭam) — a documentary}

\label{lesson:TA-C310-avanappatam}



\begin{quote}
Before the new one: say the Tamil for a cartoon, then the Tamil for a poster.
\end{quote}

\subsection*{You'll want to know: \ta{ஆவணப்படம்}}
\textbf{\ta{ஆவணப்படம்}} — \emph{āvaṇappaṭam} — \textquotedblleft{}a documentary\textquotedblright{}.

\begin{quote}
\ta{காடுகள்} \ta{பற்றிய} \ta{ஒரு} \ta{ஆவணப்படம்} \ta{பார்த்தேன்}.
\end{quote}

\begin{quote}
\emph{kāṭukaḷ paṟṟiya oru āvaṇappaṭam pārttēṉ.}
\end{quote}

\begin{quote}
I watched a documentary about forests.
\end{quote}

\ta{ஆவணம்}, a record + \ta{படம்}, a film.

\begin{cousinweb}
\textbf{\ta{ஆவணப்படம்}} (\emph{āvaṇappaṭam}) is a Tamil compound built from \ta{ஆவணம்}, \textquotedblleft{}a document\textquotedblright{}, and \ta{படம்}, \textquotedblleft{}a picture, a film\textquotedblright{}, usually traced to Sanskrit \emph{paṭa}, \textquotedblleft{}cloth\textquotedblright{} (CDIAL 7692).
\end{cousinweb}

\begin{grammarlens}[title={What you've built}]
One more word for telling what happened.
\end{grammarlens}

\subsection*{Your turn}
\begin{itemize}
\raggedright
  \item \emph{Say it:} \emph{āvaṇappaṭam}
  \item \emph{Say it:} \emph{āvaṇappaṭam}, once more
  \item \emph{Say it:} \emph{kāṭukaḷ paṟṟiya oru āvaṇappaṭam pārttēṉ.} — I watched a documentary about forests.
  \item \emph{Recall:} say the Tamil for a cartoon, then the Tamil for a poster, then say \emph{āvaṇappaṭam} again
\end{itemize}

\begin{tcolorbox}[breakable,colback=teal!4,colframe=teal!35!black,title={Before you move on}]
What does \ta{ஆவணப்படம்} mean? (\textquotedblleft{}A documentary\textquotedblright{}.) Say it once more.
\end{tcolorbox}
\section[\texorpdfstring{\ta{சித்தரி} (cittari)}{சித்தரி (cittari)}]{\ta{சித்தரி} (cittari) — to portray; to depict}

\label{lesson:TA-C310-cittari}



\begin{quote}
Before the new one: say the Tamil for a poster, then the Tamil for a documentary.
\end{quote}

\subsection*{You'll want to know: \ta{சித்தரி}}
\textbf{\ta{சித்தரி}} — \emph{cittari} — \textquotedblleft{}to portray; to depict\textquotedblright{}.

\begin{quote}
\ta{அந்தப்} \ta{படம்} \ta{கிராம} \ta{வாழ்க்கையை} \ta{அழகாகச்} \ta{சித்தரித்தது}.
\end{quote}

\begin{quote}
\emph{antap paṭam kirāma vāḻkkaiyai aḻakākac cittarittatu.}
\end{quote}

\begin{quote}
That film portrayed village life beautifully.
\end{quote}

Past (I …): \ta{சித்தரித்தேன்}, \emph{cittarittēṉ}.

What is portrayed takes -\ta{ஐ}; from \ta{சித்திரம்}, a picture.

\begin{cousinweb}
\textbf{\ta{சித்தரி}} (\emph{cittari}) is a Tamil verb built on \ta{சித்திரம்}, \textquotedblleft{}a picture\textquotedblright{}, borrowed from Sanskrit \emph{citra}, \textquotedblleft{}bright, variegated\textquotedblright{} (CDIAL 4803).
\end{cousinweb}

\begin{grammarlens}[title={What you've built}]
One more word for telling what happened.
\end{grammarlens}

\subsection*{Your turn}
\begin{itemize}
\raggedright
  \item \emph{Say it:} \emph{cittari}
  \item \emph{Say it:} \emph{cittari}, once more
  \item \emph{Say it:} \emph{antap paṭam kirāma vāḻkkaiyai aḻakākac cittarittatu.} — That film portrayed village life beautifully.
  \item \emph{Recall:} say the Tamil for a poster, then the Tamil for a documentary, then say \emph{cittari} again
\end{itemize}

\begin{tcolorbox}[breakable,colback=teal!4,colframe=teal!35!black,title={Before you move on}]
What does \ta{சித்தரி} mean? (\textquotedblleft{}To portray; to depict\textquotedblright{}.) Say it once more.
\end{tcolorbox}
\section[\texorpdfstring{\ta{ஊடகம்} (ūṭakam)}{ஊடகம் (ūṭakam)}]{\ta{ஊடகம்} (ūṭakam) — the media; a medium}

\label{lesson:TA-C310-utakam}



\begin{quote}
Before the new one: say the Tamil for a documentary, then the Tamil for to portray.
\end{quote}

\subsection*{You'll want to know: \ta{ஊடகம்}}
\textbf{\ta{ஊடகம்}} — \emph{ūṭakam} — \textquotedblleft{}the media; a medium\textquotedblright{}.

\begin{quote}
\ta{ஊடகங்கள்} \ta{அந்தச்} \ta{சம்பவத்தைப்} \ta{பெரிதாகப்} \ta{பேசின}.
\end{quote}

\begin{quote}
\emph{ūṭakaṅkaḷ antac campavattaip peritākap pēciṉa.}
\end{quote}

\begin{quote}
The media made a lot of that incident.
\end{quote}

Plural: \ta{ஊடகங்கள்}, \emph{ūṭakaṅkaḷ}.

Usually plural: \ta{ஊடகங்கள்}, the media.

\begin{cousinweb}
\textbf{\ta{ஊடகம்}} (\emph{ūṭakam}) is built on native \ta{ஊடு}, \textquotedblleft{}the middle, what comes between\textquotedblright{} (DEDR 737): the media stand between events and their readers.
\end{cousinweb}

\begin{grammarlens}[title={What you've built}]
One more word for telling what happened.
\end{grammarlens}

\subsection*{Your turn}
\begin{itemize}
\raggedright
  \item \emph{Say it:} \emph{ūṭakam}
  \item \emph{Say it:} \emph{ūṭakam}, once more
  \item \emph{Say it:} \emph{ūṭakaṅkaḷ antac campavattaip peritākap pēciṉa.} — The media made a lot of that incident.
  \item \emph{Recall:} say the Tamil for a documentary, then the Tamil for to portray, then say \emph{ūṭakam} again
\end{itemize}

\begin{tcolorbox}[breakable,colback=teal!4,colframe=teal!35!black,title={Before you move on}]
What does \ta{ஊடகம்} mean? (\textquotedblleft{}The media; a medium\textquotedblright{}.) Say it once more.
\end{tcolorbox}
